\documentclass[10pt,a4paper]{amsart}
\usepackage[margin=19mm]{geometry}
\usepackage{amsmath,amssymb,mathtools}
\usepackage[scr=rsfs,cal=euler]{mathalfa}
\usepackage{xfrac}
\usepackage[unicode,hidelinks,bookmarks=false]{hyperref}
\newcommand{\ie}{i.e.\ }
\newcommand{\cf}{cf.\ }
\newcommand{\resp}{respectively\ }
\newcommand{\Subsection}{\subsection}
\providecommand{\nobreakdash}{\nobreak\hbox{-}}
\setlength{\emergencystretch}{2em}
\multlinegap0pt
\usepackage{amssymb}
\usepackage{xcolor}
\usepackage{tikz}
\newtheorem{definition}{Definition}
\newtheorem{theorem}{Theorem}
\newcommand{\F}{\mathbb{F}}
\title{Hierarchical Quasi-cyclic Codes from Reed-Solomon and Polynomial Evaluation Codes}
\author{Emily McMillon, Kathryn Haymaker}
\date{Source: arXiv:2512.23872v1 (CC BY 4.0)}
\begin{document}
\maketitle

\noindent\emph{Source and licence.} Hierarchical Quasi-cyclic Codes from Reed-Solomon and Polynomial Evaluation Codes. Version arXiv:2512.23872v1 (CC BY 4.0), distributed by arXiv under the Creative Commons Attribution 4.0 International licence, http://creativecommons.org/licenses/by/4.0/, as recorded in the arXiv licence metadata for this exact version and shown on the abstract page https://arxiv.org/abs/2512.23872v1. Full licence notice: \texttt{licenses/arxiv-2512.23872-CC-BY-4.0.txt}.\par\smallskip

\noindent\textit{Editorial scope.} Counted unit: the article's theorem thm:qc (source0.tex lines 511--513). The article's complete original proof of it is delivered with it (source0.tex lines 515--527, one \texttt{proof} environment of 13 lines). No mathematics is written, repaired or re-derived: the statement and the proof are the article's own lines, in the article's own notation, and no retained line is substituted or re-flowed.  Retained uncounted as the source-local context the proof needs: lines 143--145.  Nothing was omitted from any retained span, so the package carries no omission record.  No retained line contains a citation, so the wrapper carries no bibliography and no bibliography entry.  No retained line contains a figure environment, an image-inclusion command or drawing code, so no image is delivered and the package declares no asset dependency.  Editorial formatting only, in addition: an AMS \texttt{amsart} preamble that restates the article's own declarations and macros used by the retained lines, namely the environments \texttt{definition}, \texttt{theorem} on separate counters, together with the standard packages those lines need.  The retained environments print in this document as Definition 1, Theorem 1 in order of appearance, whereas the article prints the same items with the numbering of its own full document.  The complete native source of the article as distributed in the arXiv e-print for this exact version is bundled unchanged as \texttt{sources/191836/source0.tex}.
\begin{definition} \label{def:pary}
    Let $\F_q$ with $q = p^r$ be a finite field of characteristic $p$ and let $\alpha$ be a primitive element of $\F_q$. Then any element of $\F_q$ can be written as $a_0 + a_1 \alpha + a_2 \alpha^2 + \dots + a_{r-1}\alpha^{r-1}$ for some collection of $a_i \in \F_p$. We will induce a \textit{$p$-ary lexicographic ordering} on these elements where $0 < 1 < \dots < p-1$ and $\alpha^0 < \alpha^1 < \dots < \alpha^{r-1}$. Similarly, for elements in $\F_q[x]$, order the coefficients of the polynomials in the $p$-ary lexicographic ordering and $x^0 < x^1< x^2 < \cdots$. For elements in $\mathbb{F}_{q}[X_1, \ldots, X_m]$,  order the coefficients of the polynomials in the $p$-ary lexicographic ordering, then use graded lexicographic ordering on the polynomials, first comparing total degree, then using lexicographic ordering on the monomial terms. 
\end{definition}

\begin{theorem} \label{thm:qc}
    Let $C$ be a $[n,k,n-k+1]_{p^r}$ Reed-Solomon code. Then $\mathcal{C}(H_C)$ is quasi-cyclic with index $p$. 
\end{theorem}

\begin{proof}
    Order the elements of $\F_{p^r}$ and the polynomials in $C$ as in Definition~\ref{def:pary}. We will show that, with this ordering, the matrix $H_C$ is quasi-cyclic with $p \times p$ circulant submatrices.

    Let $\alpha$ be a primitive element of $\F_{p^r}$. Let $\beta \in \F_q$ be arbitrary and let $q_0(x) = \sum_{i=0}^{k-1} a_i x^i$ be a polynomial in $C$ such that $a_0\in \F_{p^{r}}\setminus \F_p$.
    Then $q_0(\beta) = \sum_{i=0}^{k-1} a_i \beta^i = \sum_{i=0}^{r-1} \gamma_i \alpha^i$ for some $\gamma_i \in \F_p$. 

    Now let $q_c(x) = q_0(x) + c$ for $c \in \F_p$. Notice that 
    \begin{align*}
        q_c(\beta) &= q_0(\beta) + c = \sum_{i=0}^{r-1} \gamma_i \alpha^i + c \\
        &=  \sum_{i=1}^{r-1} \gamma_i \alpha_i + (\gamma_0+c),
    \end{align*}
    where the addition $\gamma_0 + c$ is taken modulo $p$. In particular, because of our choice in polynomial and field element ordering, this shows that the polynomial immediately after $q_0$ (modulo $p$) evaluates to the element immediately after $q_0(\beta)$, $q_0(\beta)+1$ (modulo $p$). Hence, $H_C$ is made up of $p \times p$ circulants.
\end{proof}

\end{document}
